% =============================================================================
\chapter{Applications}
\label{ch:applications}
% =============================================================================

\begin{intuition}
Ce chapitre pr\'esente les applications majeures du RL profond dans trois
domaines~: la robotique, les jeux et la finance. Pour chaque domaine, nous
discutons les d\'efis sp\'ecifiques, les architectures utilis\'ees et les
r\'esultats marquants, avec des exemples concrets d'impl\'ementation.
\end{intuition}

% -----------------------------------------------------------------------------
\section{Robotique}
% -----------------------------------------------------------------------------

\subsection{D\'efis du RL en robotique}

\begin{attention}[title=\textbf{D\'efis sp\'ecifiques \`a la robotique}]
\begin{enumerate}
    \item \textbf{S\'ecurit\'e}~: un robot physique peut se d\'etruire ou blesser
    des personnes pendant l'exploration.
    \item \textbf{Efficacit\'e d'\'echantillonnage}~: les interactions r\'eelles
    sont lentes et co\^uteuses.
    \item \textbf{Sim-to-real gap}~: les politiques apprises en simulation ne
    se transf\`erent pas directement au monde r\'eel.
    \item \textbf{Espaces continus de haute dimension}~: un robot humano\"ide
    a typiquement 20--30 degr\'es de libert\'e.
    \item \textbf{Observations bruyantes et partielles}~: capteurs imparfaits,
    latence.
\end{enumerate}
\end{attention}

\subsection{Sim-to-Real et randomisation de domaine}

\begin{definition}[Randomisation de domaine]
La \textbf{randomisation de domaine} (\emph{domain randomization}) varie
al\'eatoirement les param\`etres de la simulation (masse, friction, d\'elais,
textures) pendant l'entra\^inement, for\c{c}ant la politique \`a \^etre robuste
aux variations~:
\[
    \max_\theta \E_{\xi \sim p(\xi)}\left[J(\theta; \xi)\right]
\]
o\`u $\xi$ repr\'esente les param\`etres de simulation randomis\'es.
\end{definition}

\begin{exemple}[OpenAI --- Rubik's Cube]
En 2019, OpenAI a entra\^in\'e une main robotique Dexterous \`a r\'esoudre un
Rubik's Cube en simulation avec randomisation massive (ADR --- Automatic Domain
Randomization). La politique, entra\^in\'ee avec PPO sur des milliers de GPUs,
se transf\`ere directement au robot physique.
\end{exemple}

\subsection{Contr\^ole locomoteur}

\begin{exemple}[Locomotion avec SAC]
L'entra\^inement d'un agent quadrup\`ede dans MuJoCo avec SAC~:
\begin{itemize}
    \item \textbf{Observation}~: positions et vitesses articulaires, orientation du tronc ($\dim = 48$).
    \item \textbf{Action}~: couples articulaires ($\dim = 12$).
    \item \textbf{R\'ecompense}~: vitesse vers l'avant $- c_1 \norm{a}^2 - c_2 \norm{\text{contact}}$.
    \item R\'esultat~: d\'emarches naturelles \'emergent apr\`es $\sim 5 \times 10^6$ pas.
\end{itemize}
\end{exemple}

\begin{codeblock}[title=\textbf{Contr\^ole locomoteur avec SAC (Stable-Baselines3)}]
\begin{minted}{python}
from stable_baselines3 import SAC
from stable_baselines3.common.vec_env import SubprocVecEnv
import gymnasium as gym

def make_env(env_id):
    def _init():
        return gym.make(env_id)
    return _init

# Environnements paralleles
n_envs = 4
env = SubprocVecEnv([make_env("Ant-v5") for _ in range(n_envs)])

model = SAC(
    "MlpPolicy", env,
    learning_rate=3e-4,
    buffer_size=1_000_000,
    batch_size=256,
    tau=0.005,
    gamma=0.99,
    ent_coef="auto",
    verbose=1
)
model.learn(total_timesteps=3_000_000)
model.save("sac_ant")
\end{minted}
\end{codeblock}

\subsection{Manipulation robotique}

\begin{definition}[Hindsight Experience Replay (HER)]
\textbf{HER} (Andrychowicz et al., 2017) r\'esout le probl\`eme des r\'ecompenses
creuses en manipulation. Apr\`es un \'episode \'echou\'e, on rejoue les transitions
en rempla\c{c}ant l'objectif original par un objectif atteint~:
\[
    (s_t, a_t, r_t, s_{t+1}, g) \to (s_t, a_t, r'_t, s_{t+1}, g')
\]
o\`u $g' = s_T$ (l'\'etat final atteint) et $r'_t = R(s_t, a_t, g')$.
\end{definition}

% -----------------------------------------------------------------------------
\section{Jeux}
% -----------------------------------------------------------------------------

\subsection{Atari et DQN}

\begin{remarque}
DQN (Mnih et al., 2015) a d\'emontr\'e que le RL profond pouvait atteindre des
performances surhumaines sur 49 jeux Atari 2600 \`a partir de pixels bruts
($84 \times 84 \times 4$ frames empil\'ees). L'architecture convolutive~:
\begin{itemize}
    \item Conv $8 \times 8$, stride 4, 32 filtres + ReLU
    \item Conv $4 \times 4$, stride 2, 64 filtres + ReLU
    \item Conv $3 \times 3$, stride 1, 64 filtres + ReLU
    \item FC 512 + ReLU
    \item FC $|\mathcal{A}|$
\end{itemize}
\end{remarque}

\subsection{Go --- AlphaGo et AlphaZero}

\begin{definition}[AlphaZero]
\textbf{AlphaZero} (Silver et al., 2018) combine~:
\begin{itemize}
    \item Un r\'eseau $f_\theta(s) = (p, v)$ qui pr\'edit simultan\'ement la
    politique $p(a|s)$ et la valeur $v(s)$.
    \item Une recherche arborescente Monte Carlo (MCTS) guid\'ee par le r\'eseau.
    \item Un entra\^inement par auto-jeu~: la cible pour $p$ est la distribution
    MCTS $\boldsymbol{\pi}$~; la cible pour $v$ est le r\'esultat $z \in \{-1, +1\}$.
\end{itemize}
Perte~:
\[
    \ell(\theta) = (z - v)^2 - \boldsymbol{\pi}^\top \log \mathbf{p} + c \|\theta\|^2.
\]
\end{definition}

\subsection{Strat\'egie en temps r\'eel --- StarCraft}

\begin{exemple}[AlphaStar]
\textbf{AlphaStar} (Vinyals et al., 2019) a atteint le niveau Grandmaster \`a
StarCraft~II en combinant~:
\begin{itemize}
    \item Entra\^inement supervis\'e sur des replays humains.
    \item RL multi-agents dans une ligue de joueurs (population-based training).
    \item Architecture transformer pour g\'erer les s\'equences d'actions longues.
    \item Espace d'action combinatoire (type d'action + cible + timing).
\end{itemize}
\end{exemple}

% -----------------------------------------------------------------------------
\section{Finance}
% -----------------------------------------------------------------------------

\subsection{Gestion de portefeuille}

\begin{definition}[MDP de gestion de portefeuille]
Le probl\`eme de gestion de portefeuille se formalise comme~:
\begin{itemize}
    \item \textbf{\'Etat}~: prix historiques, indicateurs techniques, positions
    courantes, solde.
    \item \textbf{Action}~: vecteur d'allocation $\mathbf{w} \in \Delta^K$ o\`u
    $w_k$ est la fraction allou\'ee \`a l'actif $k$.
    \item \textbf{R\'ecompense}~: rendement logarithmique ajust\'e du risque~:
    \[
        R_t = \log\left(\sum_k w_k \frac{p_k^{t+1}}{p_k^t}\right) - c_{\text{transaction}}
        - \lambda_{\text{risk}} \cdot \text{CVaR}_\alpha
    \]
\end{itemize}
\end{definition}

\begin{codeblock}[title=\textbf{Environnement de trading simplifi\'e}]
\begin{minted}{python}
import numpy as np
import gymnasium as gym
from gymnasium import spaces

class TradingEnv(gym.Env):
    def __init__(self, prices, window=20, initial_balance=10000):
        super().__init__()
        self.prices = prices  # (T, n_assets)
        self.window = window
        self.n_assets = prices.shape[1]
        self.initial_balance = initial_balance
        self.observation_space = spaces.Box(
            low=-np.inf, high=np.inf,
            shape=(window, self.n_assets + 1))  # prix + position
        self.action_space = spaces.Box(
            low=0, high=1, shape=(self.n_assets,))
        self.reset()

    def reset(self, seed=None, options=None):
        super().reset(seed=seed)
        self.t = self.window
        self.balance = self.initial_balance
        self.holdings = np.zeros(self.n_assets)
        return self._get_obs(), {}

    def _get_obs(self):
        price_window = self.prices[self.t - self.window:self.t]
        # Rendements normalises
        returns = np.diff(price_window, axis=0) / price_window[:-1]
        position = np.tile(self.holdings / (self.balance + 1e-8),
                          (self.window, 1))
        return np.concatenate([returns, position[:returns.shape[0]]], axis=1)

    def step(self, action):
        weights = action / (action.sum() + 1e-8)
        # Rendements des actifs
        current_prices = self.prices[self.t]
        next_prices = self.prices[self.t + 1]
        asset_returns = next_prices / current_prices - 1

        portfolio_return = np.dot(weights, asset_returns)
        transaction_cost = 0.001 * np.sum(
            np.abs(weights - self.holdings / (self.balance + 1e-8)))

        self.balance *= (1 + portfolio_return - transaction_cost)
        self.holdings = weights * self.balance
        self.t += 1

        reward = np.log(1 + portfolio_return - transaction_cost)
        terminated = self.t >= len(self.prices) - 1
        return self._get_obs(), reward, terminated, False, {
            "balance": self.balance}
\end{minted}
\end{codeblock}

\subsection{Ex\'ecution optimale d'ordres}

\begin{definition}[Probl\`eme d'ex\'ecution optimale]
Un trader doit ex\'ecuter un ordre de $Q$ parts en $T$ p\'eriodes en minimisant
l'impact de march\'e~:
\begin{itemize}
    \item \textbf{\'Etat}~: $s_t = (q_t, t, \text{features de march\'e})$ o\`u
    $q_t$ est la quantit\'e restante.
    \item \textbf{Action}~: $a_t$ quantit\'e \`a ex\'ecuter au temps $t$.
    \item \textbf{Contrainte}~: $\sum_t a_t = Q$ et $a_t \geq 0$.
    \item \textbf{R\'ecompense}~: $-a_t \cdot \text{slippage}(a_t)$.
\end{itemize}
\end{definition}

% -----------------------------------------------------------------------------
\section{Autres applications}
% -----------------------------------------------------------------------------

\subsection{RLHF --- Alignement de mod\`eles de langage}

\begin{definition}[RLHF]
Le \textbf{Reinforcement Learning from Human Feedback} (Ouyang et al., 2022)
aligne les LLM avec les pr\'ef\'erences humaines~:
\begin{enumerate}
    \item Entra\^iner un mod\`ele de r\'ecompense $R_\psi(x, y)$ sur des
    pr\'ef\'erences humaines $(y_w \succ y_l | x)$.
    \item Optimiser le LLM $\pi_\theta$ par PPO~:
    \[
        \max_\theta \E_{x \sim \mathcal{D}, y \sim \pi_\theta}\left[
        R_\psi(x, y) - \beta \, \text{KL}[\pi_\theta \| \pi_{\text{ref}}]
        \right].
    \]
\end{enumerate}
\end{definition}

\subsection{D\'ecouverte de m\'edicaments}

\begin{exemple}[G\'en\'eration mol\'eculaire]
Le RL est utilis\'e pour g\'en\'erer des mol\'ecules optimisant des propri\'et\'es
pharmacologiques (solubilit\'e, affinit\'e, toxicit\'e). L'agent construit
une mol\'ecule atome par atome ou fragment par fragment~:
\begin{itemize}
    \item \textbf{\'Etat}~: graphe mol\'eculaire partiel.
    \item \textbf{Action}~: ajouter un atome/liaison.
    \item \textbf{R\'ecompense}~: score de \emph{docking}, QED, SA score.
\end{itemize}
\end{exemple}

\subsection{Contr\^ole de plasma de fusion}

\begin{exemple}[DeepMind \& EPFL (2022)]
Le RL a \'et\'e utilis\'e pour contr\^oler la forme du plasma dans le tokamak
TCV~:
\begin{itemize}
    \item \textbf{\'Etat}~: mesures magn\'etiques (200+ capteurs).
    \item \textbf{Action}~: tensions des bobines magn\'etiques (19 actionneurs).
    \item \textbf{R\'ecompense}~: \'ecart \`a la forme cible du plasma.
    \item L'agent PPO contr\^ole avec succ\`es des configurations de plasma
    jamais tent\'ees auparavant.
\end{itemize}
\end{exemple}

% -----------------------------------------------------------------------------
\section{Bonnes pratiques}
% -----------------------------------------------------------------------------

\begin{formules}[title=\textbf{Guide de choix d'algorithme}]
\begin{itemize}
    \item \textbf{Actions discr\`etes, observation simple}~: DQN, Rainbow, PPO.
    \item \textbf{Actions continues}~: SAC, TD3, PPO.
    \item \textbf{Efficacit\'e d'\'echantillonnage critique}~: SAC, TD3 (off-policy).
    \item \textbf{Stabilit\'e et simplicit\'e}~: PPO.
    \item \textbf{Multi-agents coop\'eratifs}~: MAPPO, QMIX.
    \item \textbf{Contraintes de s\'ecurit\'e}~: PPO-Lagrangien, CPO.
\end{itemize}
\end{formules}

\begin{attention}[title=\textbf{Pi\`eges courants}]
\begin{enumerate}
    \item Ne pas normaliser les observations et les r\'ecompenses.
    \item Utiliser un taux d'apprentissage trop \'elev\'e.
    \item Replay buffer trop petit pour les m\'ethodes off-policy.
    \item Ignorer les graines al\'eatoires~: toujours rapporter les r\'esultats
    sur 5--10 graines avec intervalles de confiance.
    \item Ne pas v\'erifier l'impl\'ementation sur un probl\`eme simple d'abord.
\end{enumerate}
\end{attention}

% -----------------------------------------------------------------------------
\section{Exercices}
% -----------------------------------------------------------------------------

\begin{exercice}[Sim-to-real]
Entra\^inez un agent PPO sur \texttt{Pendulum-v1} avec randomisation de la masse
et de la friction. \'Evaluez la robustesse en testant sur des param\`etres non
vus pendant l'entra\^inement.
\end{exercice}

\begin{exercice}[Trading]
Impl\'ementez l'environnement de trading ci-dessus avec des donn\'ees historiques
r\'eelles (par ex.\ Yahoo Finance). Entra\^inez SAC et comparez avec une
strat\'egie buy-and-hold et une strat\'egie de moyenne mobile.
\end{exercice}

\begin{exercice}[MCTS + RL]
Impl\'ementez une version simplifi\'ee d'AlphaZero pour le jeu de Connect Four.
Combinez un r\'eseau politique-valeur avec MCTS et entra\^inez par auto-jeu.
\end{exercice}

\begin{exercice}[HER]
Impl\'ementez HER pour un environnement de manipulation
\texttt{FetchReach-v3}. Comparez les courbes d'apprentissage avec et sans HER.
\end{exercice}
